Benchmarks of inferred gene regulatory networks rely on incomplete, mutually dependent, unevenly curated references.
We formulate multi-reference evaluation as Bayesian latent-class estimation: \nNAnnotHeldin{} public transcription factor (TF)--target resources, grouped by evidence modality, are imperfect annotators of a latent true edge, with family-specific dependence and hierarchical pooling (Family-HB), fitted by posterior-mode search and Hamiltonian Monte Carlo.
Family-HB is no more robust to resource dropout than a Dawid--Skene model and does not outperform simple combination rules on held-out ChIP-seq references and \nCohN{} CRISPR screens (\nRepPassAllD{} studies evaluable).
Its posterior is bimodal, and in tissue universes literature consensus and copying are only weakly distinguishable.
Post hoc analyses show that coverage is not missing at random, which an informative-observation model does not repair.
Precision gains in context-matched universes reflect prevalence.
Label-free claim tiers controlling the posterior expected false discovery rate are nonetheless enriched for held-out binding (\nGovOeEnetMin{}- to \nGovOeEnetMax{}-fold on a specific ChIP network, TF-matched) and show, in \nRepPassPowD{} adequately powered screens (global universe), a modest TF-specific association with perturbation response.
For scGPT and Geneformer attention in three tissues, almost no top-ranked scGPT edge qualifies as a supported claim, and scGPT attention largely reflects the number of genes expressed per cell.
